\chapter*{Abstract}
\addcontentsline{toc}{chapter}{Abstract}

Unmanned Aerial Vehicles equipped with visible and thermal cameras can support
wilderness Search and Rescue, but imperfect sensor registration, missing or
uninformative modalities, limited data, and training variability complicate
reliable detector comparison. This thesis studies RGB--thermal person detection
on the WiSARD dataset, starting from a dual-backbone RT-DETR detector and an
inherited Feature Alignment Module (FAM). It contributes a reliability-aware
evaluation and extension of this starting point through paired multi-seed
comparisons, internal diagnostics, a complete-video validation protocol,
matched full-data confirmation, and transfer to other detector families.

Principal controlled comparisons use five paired seeds and fixed training
budgets without early stopping. In the validation-led protocol, Stage A selects
\texttt{best} on a held-out complete video and retains the final checkpoint as
a diagnostic. A candidate satisfying the aggregate promotion rule is retrained
from the pretrained initialization on all training data in Stage B and compared
with a freshly trained matched control using final checkpoints only.

In the historical locked comparison, standard FAM improves Base RT-DETR in all
five seeds, increasing mean \mapfifty{} from \(0.3080\pm0.0677\) to
\(0.3780\pm0.0440\). Its main behavioral effect is higher recall: at confidence
0.25, the mean paired increase is 0.0990 and occurs in every seed. Internal
inspection confirms learned transformations at P3 and P4 without establishing
physical RGB--IR registration, and reveals a compensated P5 collapse in one
high-performing checkpoint. The validation-led campaign rejects the tested P2,
global \(800\times800\), post-FAM gating, scalar residual, mixed-consistency,
and box-guided variants. Reliability-Conditioned Residual Alignment (RCRA) is
the only candidate promoted from Stage A. In matched full-data Stage B, it
obtains \(0.3753\pm0.0538\), against \(0.3569\pm0.0422\) for the matched FAM
control, but wins only three of five seeds and therefore does not replace FAM.

A post-selection audit of the matched Stage-B FAM checkpoints on the same 708
paired frames shows that fusion exceeds VIS in all five seeds by
\(0.0373\pm0.0223\) \mapfifty{}. Masked IR performs poorly against VIS
annotations, whereas native-IR evaluation confirms that the thermal branch
remains functional; the two conditions use different coordinate contracts.
On two previously unused internal WiSARD acquisitions, historical FAM exceeds
historical Base in all ten acquisition--seed comparisons, with a macro-average
gain of 0.1257. Controlled IR translations and rescaling nevertheless do not
establish a universally more favorable signed relative-drop curve for FAM.

Architectural transfer is also conditional. A separate five-seed, 200-epoch
YOLOv10 protocol finds no consistent FAM benefit. RT-DETRv2 FAM has a positive
mean Stage-A gain of 0.0220 but only three positive pairs, while YOLO26 FAM
reduces the best-validation mean by 0.0234 and wins only one pair. These are
within-family comparisons rather than a detector leaderboard. Standard FAM
therefore remains the selected RT-DETR configuration, while the broader result
is methodological: multimodal improvements require paired replication,
explicit coordinate contracts, and predeclared model-selection rules.

\vspace{1em}
\noindent\textbf{Keywords:} RGB--thermal fusion; Search and Rescue; UAV;
RT-DETR; RT-DETRv2; YOLOv10; YOLO26; missing modalities; feature alignment;
reproducibility.
